\section{Uso de la ciencia de datos en la inclusión financiera del sector agro}\label{sec:inclusion}

\textbf{De un rendimiento por zona a uno por parcela.} La capa de crédito sustituye el rendimiento esperado del
paquete tecnológico por la predicción de cada parcela con su incertidumbre. Calcula la probabilidad de no cubrir el crédito o el
costo total. La probabilidad combina los errores del modelo con
la fracción de superficie que el municipio pierde históricamente, porque si la cosecha se pierde el ingreso es cero. Un
semáforo la resume en verde bajo 10\,\%, amarillo de 10 a 30\,\% y rojo desde 30\,\%. Con el crédito, 23 de las 59
parcelas salen en verde, 30 en amarillo y 6 en rojo. Con el costo total salen 55 en rojo, porque el paquete incluye
conceptos no financiados, como la renta.

\textbf{A quién sirve.} La comercializadora que compra el grano puede priorizar visitas de campo. El analista ordena
parcelas por riesgo y el productor ve cuánto debe cosechar para pagar.

\textbf{Alcance.} La probabilidad describe la incertidumbre de un año parecido a 2025 y no incluye sequías, heladas ni
cambios de precio. Por eso no es una probabilidad de incumplimiento ni sirve para tarificar seguros, que además
enfrentan riesgo base~\cite{jensen2017}.

\textbf{Momento de uso.} El modelo usa rendimientos de vecinas de la misma temporada, que existen solo después de su
cosecha. Es una estimación de cierre de temporada. Para usarlo antes de cosechar harían falta rendimientos de ciclos
anteriores, que el reto no incluye.

\textbf{Escala municipal.} Con el SIAP de 2003 a 2025 ajustamos un modelo municipal~\cite{siap2026}. En
municipios a más de 50\,km de los de entrenamiento, la temperatura máxima de julio y agosto más kriging de residuos
obtuvo un RMSE de 0.44\,t/ha contra 0.57 del kriging solo. Al repetir de forma anidada
la elección de esa variable, la ventaja baja a unos 0.07\,t/ha.
